% Response to reviewers -- ACML 2026 collection, Machine Learning (Springer)
% Submission 62d20e3d-c3eb-4568-b964-7b09cf6f583b
%
% Standalone: compiles with pdflatex on a plain TeX installation and does not need the
% Springer class. Every reviewer point gets the same four parts -- the reviewer's words, what we
% did, where in the manuscript it lives, and the revised text quoted verbatim -- so a reviewer can
% check any single point without reading the rest.
\documentclass[11pt,a4paper]{article}

\usepackage[T1]{fontenc}
\usepackage[utf8]{inputenc}
\usepackage[margin=2.4cm]{geometry}
\usepackage{amsmath}
\usepackage{booktabs}
\usepackage{longtable}
\usepackage{enumitem}
\usepackage{xcolor}
\usepackage[colorlinks=true,linkcolor=linkblue,urlcolor=linkblue,citecolor=linkblue]{hyperref}

\definecolor{linkblue}{HTML}{0F4D92}
\definecolor{revgrey}{HTML}{4D4D4D}
\definecolor{quoterule}{HTML}{0F4D92}
\definecolor{rulegrey}{HTML}{C4C4C4}

\setlength{\parindent}{0pt}
\setlength{\parskip}{0.55em}

% The reviewer's own words.
\newenvironment{reviewer}
  {\begin{list}{}{\leftmargin=1.1em\rightmargin=0.6em}\item[]\itshape\color{revgrey}\small}
  {\end{list}}

% Text quoted verbatim from the revised manuscript, marked with a rule so it cannot be mistaken
% for our commentary in this letter.
\newenvironment{revised}
  {\begin{list}{}{\leftmargin=1.1em\rightmargin=0.6em}\item[]%
   \color{black}\small\hspace{-1.1em}\textcolor{quoterule}{\rule[-0.4ex]{2pt}{1.05em}}\hspace{0.5em}\ignorespaces}
  {\end{list}}

\newcommand{\where}[1]{{\textbf{Where addressed:} #1\par}}
\newcommand{\rpoint}[2]{\subsection*{#1 --- #2}}

\title{\vspace{-1.6cm}\Large\bfseries Response to Reviewers}
\date{}
\author{}

\begin{document}
\maketitle
\vspace{-2.2cm}

\noindent\textcolor{rulegrey}{\rule{\textwidth}{1pt}}

\textbf{Manuscript.} Worst-Group Conformal Reliability under Spurious Correlation Is More a
Calibration Problem than a Representation Problem: A Confound-Controlled Multi-Backbone Study

\textbf{Submission ID.} 62d20e3d-c3eb-4568-b964-7b09cf6f583b \quad ACML 2026 collection,
\emph{Machine Learning}

\textbf{Authors.} One Octadion, Novanto Yudistira, Yuta Nakashima

\noindent\textcolor{rulegrey}{\rule{\textwidth}{1pt}}

We thank the three reviewers and the editor for reviews that were specific enough to act on. Two of
the comments identified genuine errors in our analysis, and acting on them changed several reported
numbers; we say so explicitly below rather than presenting the revision as a clarification. All
reviewer requests have been implemented, including the two that required new experiments.

The revision rests on new runs: a four-backbone grid ($36{,}000$ records), a four-backbone
calibration ablation across three calibration policies ($64{,}800$), a four-backbone
predicted-group study ($10{,}800$), and an end-to-end fine-tuning study in which the representation
itself is retrained ($4{,}320$). Every table in the paper has been recomputed from these.

\section*{Summary of the principal changes}

\begin{longtable}{@{}p{0.5em}p{11.4cm}p{2.9cm}@{}}
\toprule
\textbf{\#} & \textbf{Change} & \textbf{Prompted by} \\
\midrule
\endfirsthead
\toprule
\textbf{\#} & \textbf{Change} & \textbf{Prompted by} \\
\midrule
\endhead
1 & \textbf{New experiment:} the backbone is fine-tuned end-to-end under three objectives, so the
representation itself moves rather than only the head. New Section~5.2 and Table~4. & R1.2, R2.1 \\
\addlinespace
2 & \textbf{Backbones extended from two to four}, spanning two architecture families and three
pretraining regimes. Every results table now reports eight backbone--dataset cells. & R1.3 \\
\addlinespace
3 & \textbf{Title narrowed} from an absolute dichotomy to a comparative claim, and the same
softening applied in fifteen further places. & R1.2, R2.1, R3 \\
\addlinespace
4 & \textbf{Proposition~1(ii) rewritten} around an exact identity and a bound that holds, with an
explicit remark on what the divergences do and do not give. & R2.2 \\
\addlinespace
5 & \textbf{Uncertainty quantification replaced} with a two-stage cluster bootstrap that respects
calibration splits nested within trained models. & R2.3 \\
\addlinespace
6 & \textbf{Four claims withdrawn} and two matched-divergence values retracted, because the
corrected analysis does not support them. & R1.3, R2.3, R2.6, R3 \\
\addlinespace
7 & \textbf{Sub-target coverage explained with a direct measurement:} mean-over-groups coverage is
$0.9024$ against a nominal $0.900$. & R1.1, R3 \\
\addlinespace
8 & \textbf{Sensitivity analysis over the excluded runs}, and the exclusion rules described as
pre-specified rather than preregistered. & R2.4 \\
\addlinespace
9 & \textbf{Predicted-group mechanics, supervision, error propagation and ethics} written out in
full; the $\rho$-sweep claim scoped to group-prior shift. & R2.5 \\
\addlinespace
10 & \textbf{Reproducibility appendix expanded}; literature comparison table added; requested
citation added. & R1.3, R1.5, R2.6 \\
\bottomrule
\end{longtable}

The revised manuscript runs to about 23 pages against the 20 that applied at submission. We would
rather report the requested experiments and analyses in full than compress them, and we are glad to
shorten if the editorial office requires it.

\section*{Reviewer 1}

\rpoint{R1.1}{the target is $0.90$ but the empirical values are $0.84$--$0.89$}

\begin{reviewer}
``In Section 3 \ldots the authors state that $(1-\alpha=0.90)$ is the target throughout. However,
their empirical values are around $0.84$--$0.89$, which are below the stated $0.90$ target. The
authors should clearly explain the reason for this difference.''
\end{reviewer}

The reviewer is right that the paper reported a sub-target number without an adequate explanation.
The cause is that we report a \textbf{minimum over four noisy per-group coverages}, which is a
downward-biased statistic even when every group is individually valid. We now demonstrate this with
a direct measurement rather than an argument.

\where{Section~5.1 (end), Section~5.7, and Appendix~B (including a new figure).}

\begin{revised}
Mondrian's level settles at $0.854$--$0.886$, below the nominal $0.90$. That is a property of the
statistic rather than a validity failure, and we quantify it directly: the unweighted mean over
groups, which is what group-conditional calibration actually targets, is $\mathbf{0.9024}$ across
the same eight cells (range $0.8995$--$0.9056$), while the \emph{minimum} over four noisy per-group
coverages sits $0.018$--$0.045$ below it.
\end{revised}

Appendix~B additionally simulates an \emph{exactly valid} Mondrian procedure with no model and no
data, and a new figure puts our runs against that law directly:

\begin{revised}
The blue line is where an \emph{exactly valid} Mondrian procedure puts that minimum, simulated at
our own per-group calibration and test counts with no model and no data: $\mathbf{0.8695}$ against
an observed $\mathbf{0.8660}$ on Waterbirds and $\mathbf{0.8847}$ against $\mathbf{0.8800}$ on
CelebA, with $98\%$ and $95\%$ of runs inside its central $95\%$ range.
\end{revised}

The two effects are worth separating. The closed-form $1.03\,\sigma$ figure covers only the
min-over-$k$ selection effect and predicts $0.013$--$0.030$ against an observed $0.018$--$0.045$,
so on its own it accounts for the level's direction and order of magnitude but not every point of
it. Simulating the full Mondrian law additionally captures the second effect --- each threshold is
estimated from that group's own small calibration sample --- and the residual is then
$0.004$--$0.005$ of coverage. We report that residual rather than claim an exact match.

\rpoint{R1.2}{the backbones are frozen, so why the title claim?}

\begin{reviewer}
``The authors freeze the neural network backbones \ldots Therefore, why do the authors claim in the
title \ldots? The authors should clarify this claim or provide additional experiments where the
representation itself is trained or fine-tuned.''
\end{reviewer}

We have done both. This was the most substantial change to the paper.

\where{New Section~5.2 and Table~4 (experiment); the title, abstract, Section~5.1 heading,
discussion, conclusion and related work (narrowing); Section~3 (terminology).}

We fine-tune the backbone \textbf{end-to-end} under three objectives --- plain ERM, GroupDRO and
group reweighting --- with five fine-tuning seeds each, holding the head fixed at plain ERM so the
representation is the only thing that moves. The section opens with a \textbf{manipulation check},
because a representation-level experiment answers nothing unless the intervention actually moved
the representation:

\begin{revised}
GroupDRO fine-tuning raises worst-group accuracy from $0.516$ to $0.738$ on Waterbirds ($+0.222$)
and from $0.396$ to $0.470$ on CelebA ($+0.074$): the lever was pulled, and by the largest margin
any intervention in this paper produces. Reweighting did not help ($0.466$ and $0.364$, both
slightly below ERM). We report it regardless --- a representation that fails to improve is still a
representation that changed, and dropping it would be selecting on the outcome.
\end{revised}

The dissociation survives. Under marginal calibration the three representations span $0.249$ of
worst-group coverage on Waterbirds; under Mondrian they sit within $0.004$ of one another.

\textbf{Title narrowed.} ``\ldots Is a Calibration Problem, \emph{Not} a Representation Problem''
has become ``\ldots Is \emph{More} a Calibration Problem \emph{than} a Representation Problem''. We
traced the same absolute claim through the rest of the paper and softened it in fifteen further
places, since softening only the title would not have answered the objection.

We also corrected a terminological conflation the reviewer's question exposed: the submitted paper
used ``representation'' for the frozen backbone, which it never varied. Throughout the revision,
\textbf{backbone} means the frozen feature extractor, \textbf{representation} means what the
fine-tuning study changes, and \textbf{head} means the last-layer classifier.

\rpoint{R1.3}{only two backbones; compare more clearly with related work}

\begin{reviewer}
``Only two backbones are used, which limits the generalizability of the results. In addition, the
experimental results should be compared more clearly with those of other related studies.''
\end{reviewer}

\where{Section~4 and every results table (backbones); new Table~2 and three paragraphs before
Section~5 (literature comparison).}

\textbf{Backbones: two $\rightarrow$ four,} chosen to span two architecture families and three
pretraining regimes --- an ERM-trained ResNet-50 ($d=2048$), CLIP ViT-B/32 ($512$, image--text
contrastive), DINOv2 ViT-B/14 ($768$, self-supervised, never label-trained), and a supervised
ImageNet ViT-B/16 ($768$). Every results table now reports all eight backbone--dataset cells. The
substantive finding is that the two added backbones behave like the two originally reported: the
Mondrian cross-training spread never exceeds $0.024$ in any cell.

\textbf{Literature comparison.} The comparison is unflattering and we present it as such --- our
ResNet-50 arms are $10$--$24$ points below published worst-group accuracies on Waterbirds --- and
we state the three protocol differences that explain it: our arms refit only the last layer on
cached features (where the published GroupDRO and AFR train the network), our ResNet-50 uses a
deliberately light recipe and a $30{,}000$-image CelebA training subsample, and worst-group
accuracy is read on our pooled evaluation domain rather than the benchmark test split. DFR, the one
arm whose protocol matches its published counterpart, is correspondingly the closest ($0.851$
against $0.883$ on CelebA).

We also explain why this does not bear on the paper's claims: every result is a
\textbf{within-cell} contrast computed on the same posteriors, so a weaker backbone lowers all arms
in a cell together and leaves the calibration contrast intact. The four backbones span a
$0.20$--$0.95$ range of worst-group accuracy and the dissociation is unchanged across all of them.

\rpoint{R1.4}{more independent training seeds}

\begin{reviewer}
``The authors should increase the number of independent training seeds to ensure the reproducibility
and stability of the reported results.''
\end{reviewer}

\where{Section~4, new paragraph; Appendix~D (seeds and splits).}

The main grid and the end-to-end fine-tuning study now run \textbf{five} training seeds; the
calibration-policy ablation runs three, and the manuscript states the design as $\ge 3$ training
seeds $\times \ge 10$ calibration splits rather than quoting a single figure that would not hold
everywhere. In raising the count we found something we must disclose, because presenting it as
stability would be misleading: \textbf{two of the five arms are deterministic.} ERM and AFR fit with
L-BFGS, whose solver ignores the random seed, so their across-seed standard deviation is exactly
$0.000$ and their base accuracy is identical to five decimals in every cell.

\begin{revised}
Running more seeds therefore adds replicates for DFR, balanced subsampling and GroupDRO-LL
(across-seed SD $0.003$--$0.027$ in worst-group accuracy) but not for ERM or AFR. We report this
rather than present zero variance as stability.
\end{revised}

This is also why every interval in the revision is a cluster bootstrap resampling seeds and then
splits within them: the calibration splits inside a seed share one trained model.

\rpoint{R1.5}{add the requested reference}

\begin{reviewer}
``The authors must add the following recent study to the references: `New unfreezing strategy of
transfer learning in satellite imagery\ldots', Scientific African, 2024.''
\end{reviewer}

\where{Section~5.2, in the fine-tuning protocol.}

Added, and placed where it is substantively relevant rather than appended to a list. Because the
revision now fine-tunes backbones end-to-end, how much of a pretrained network to unfreeze is a
choice we actually make:

\begin{revised}
How much of a pretrained network to unfreeze is itself consequential, and staged or partial
unfreezing schedules change what a fine-tuned representation encodes [citation]; we unfreeze the
whole network, so the representation lever is pulled as far as each objective allows and this is
the strongest test of whether calibration still dominates that our setting admits.
\end{revised}

\section*{Reviewer 2}

\rpoint{R2.1}{scope of the main claim}

\begin{reviewer}
``All training interventions modify only the last-layer classifier while the backbone remains
frozen. Please add a representation-changing robustness experiment if feasible; otherwise, narrow
the title, abstract, and conclusions\ldots''
\end{reviewer}

\where{New Section~5.2 (experiment); title, abstract (fourth sentence), and conclusions
(narrowing).}

We have done both, as described under R1.2: the new Section~5.2 retrains the representation
end-to-end, and the title, abstract and conclusions have been narrowed regardless. The claim is now
comparative throughout, and the abstract states the scope in its fourth sentence --- ``We test that
intuition on \emph{both} levers'' --- so a reader encounters it before the results.

\rpoint{R2.2}{Proposition~1(ii)}

\begin{reviewer}
``Wasserstein-1 or KS divergence alone generally does not imply a monotonic reduction in coverage at
the pooled threshold. Please state sufficient assumptions for such a relationship or present the
divergence--coverage link as an empirical hypothesis.''
\end{reviewer}

The reviewer is correct and the earlier claim did not follow. Proposition~1(ii) is now built on an
\textbf{exact identity} rather than an assertion about divergences:
\[
(1-\alpha) - F_{\mathrm{wg}}(\hat q)
= (1-\pi)\bigl[F_{\neg\mathrm{wg}}(\hat q) - F_{\mathrm{wg}}(\hat q)\bigr],
\qquad\text{hence}\qquad
(1-\alpha) - F_{\mathrm{wg}}(\hat q) \le (1-\pi)\,D_{\mathrm{KS}} .
\]

\where{Section~3, Proposition~1(ii) and the remark that follows it.}

\begin{revised}
$W_1$ admits no such bound --- it integrates the gap over the line while the shortfall depends on it
at one point, so the gap can concentrate on a vanishing interval with $W_1 \to 0$ and the shortfall
fixed --- so we use $W_1$ only as a descriptive summary of heterogeneity. Neither divergence
\emph{orders} coverage without a direction assumption: $P_{\mathrm{wg}}$ stochastically \emph{below}
the pool has the same $W_1$ and $D_{\mathrm{KS}}$ as its mirror image above it but yields a coverage
surplus. Monotonicity needs $P_{\mathrm{wg}} \succeq_{\mathrm{st}} P_{\neg\mathrm{wg}}$, which holds
here by construction, the worst group being defined as the lowest-coverage one --- an assumption
about which group is worst, not a consequence of the divergence.
\end{revised}

All four statements were verified numerically before being written.

\rpoint{R2.3}{the accuracy-matched analysis}

\begin{reviewer}
``The current interpolation uses relatively few independent training seeds, and several robust
models show no overlap in accuracy with ERM. More controlled overlap and independent runs would
strengthen the analysis; otherwise, limit the conclusion to matched settings and use a hierarchical
uncertainty analysis that respects repeated calibration splits within the same trained model.''
\end{reviewer}

This comment led to the largest correction in the revision, and the reviewer's suspicion was better
founded than we realised.

\where{Section~5.5 (rewritten), Table~3 (recomputed), new Figure~4, and Appendix~D.}

\textbf{What we found.} The submitted analysis pooled each arm's (accuracy, divergence) points
across seeds \emph{and} calibration splits. Measured separately, ERM's per-seed base accuracy is
identical to five decimals in every cell --- its solver is deterministic --- while its across-split
range is $0.02$--$0.03$. So the entire accuracy axis ERM contributed was \emph{which evaluation rows
were drawn}, not \emph{which model was trained}. Matching on it compares ERM on unlucky draws
against a robust arm on lucky ones. The accompanying bootstrap over pooled points also treated
nested rows as independent draws, which is exactly the hierarchical-uncertainty problem the reviewer
raises.

\textbf{What we did.} We recomputed at model level (one accuracy per training seed) with a two-stage
cluster bootstrap that resamples seeds and then splits within them. ERM's accuracy support is then a
single \textbf{point}, and overlapping support exists in \textbf{one of the eight cells}:
GroupDRO-LL on ResNet-50/Waterbirds, whose per-seed accuracies bracket ERM's $0.9399$. There the
reduction is real and \textbf{replicates} the earlier estimate ---
$\Delta_{\mathrm{APS}} = +0.084\,[+0.076, +0.095]$,
$\Delta_{\mathrm{RAPS}} = +0.084\,[+0.074, +0.094]$,
$\Delta_{\mathrm{THR}} = +0.150\,[+0.146, +0.158]$.

\textbf{What we withdraw.} The two other matched values in the submitted Table~3 --- GroupDRO-LL on
CLIP/Waterbirds ($+0.100$) and balanced subsampling on ResNet-50/Waterbirds ($+0.059$) --- rested on
the split-level spread. At model level neither arm's accuracy support reaches ERM's, and both are
now reported as unmatched.

Following the reviewer's instruction, we \textbf{limit H3 to that one matched setting} and draw no
accuracy-matched conclusion elsewhere. We also note that the reason the comparison is unavailable is
itself the more robust reading: on CelebA the robust arms raise base accuracy to $0.81$--$0.91$
against ERM's $0.68$--$0.70$, so no common operating point exists at all.

\rpoint{R2.4}{impact of the excluded runs, and preregistration status}

\begin{reviewer}
``The failed AFR runs on CelebA and the excluded GroupDRO seed are informative for comparing
training interventions. Please include them in a sensitivity analysis and clarify whether the
exclusion rules were preregistered or only pre-specified.''
\end{reviewer}

\where{Section~4 (heading and wording), Appendix~C (sensitivity analysis), Appendix~C (AFR
diagnosis).}

\textbf{Status of the rules.} They were \textbf{pre-specified, not preregistered.} The per-method
accuracy floors were fixed in the codebase on 11 June 2026, three months before any run reported
here, so they are pre-specified with respect to every number in the paper --- but they were never
lodged in a public registry, and we do not claim otherwise. We have also corrected the paper's own
heading, which read ``Pre-registered hypotheses'' and now reads ``Pre-specified hypotheses'' with
the distinction stated.

\textbf{Sensitivity analysis.} Every excluded arm was evaluated and its records retained, so this
needed no re-run. Including them, the dissociation is unaffected: the Mondrian cross-training spread
rises only to $0.028$ against $0.024$. \textbf{One thing does change, and we report it against
ourselves:} the pre-specified C1 criterion --- Mondrian within $0.02$ of target \emph{and} marginal
significantly below target for \emph{every} method --- is met in exactly one of eight cells, and
that single cell \textbf{flips to failing} when the gated AFR arm is included. We therefore rest the
claim on the spread magnitudes rather than on the pass/fail flag.

\textbf{AFR's collapse is not an untuned hyperparameter.} Since excluding a method on a
hyperparameter we chose would be unfair, we investigated. The mechanism is a class-prior inversion:
at the default $\gamma=2$ the reweighting drives the head to predict the positive class for
$\approx 86\%$ of examples where the true rate is $\approx 15\%$, so a \emph{majority} group becomes
the worst group. We then computed an \textbf{oracle upper bound} --- the best $\gamma$ chosen with
knowledge of the evaluation set, which no selection rule could beat --- of $0.327$, $0.427$, $0.445$
and $0.477$ across the four CelebA backbones, all far below the $0.75$ floor. AFR's exclusion is a
property of the method in this regime, not of our choice of $\gamma$.

\rpoint{R2.5}{shift and predicted-group deployment claims}

\begin{reviewer}
``The rho sweep mainly appears to change group mixture/correlation strength rather than the
within-group data distribution, so the robustness claim should be scoped accordingly \ldots For
predicted-group Mondrian, describe exactly how predicted attributes determine thresholds, what
supervision is required, how probe errors affect coverage, and discuss the privacy or ethical
implications of inferring the CelebA \emph{Male} attribute.''
\end{reviewer}

All five requests implemented.

\where{Section~5.7 ($\rho$-sweep scoping); Section~5.8 (mechanics, supervision, probe error);
Section~6 and the ethics statement (privacy).}

\textbf{$\rho$-sweep scoping.} The reviewer's reading is correct.

\begin{revised}
The $\rho$ sweep changes the \emph{group mixture} of the evaluation distribution \ldots while the
class-conditional feature distribution within each of the four groups is untouched. It is therefore
a group-prior shift, a special case of label-and-attribute shift, and \emph{not} a general covariate
shift \ldots Our stability claim is scoped accordingly.
\end{revised}

\textbf{Mechanics.} The label is never predicted; only the attribute is. The stratum is
$\hat g = 2y + \hat a$; calibration points are binned by $\hat g$, each bin's
$\lceil(1-\alpha)(n+1)\rceil$-th smallest score becomes its threshold, and a test point is admitted
if its score is at most its own bin's threshold. Coverage is always tallied against the
\textbf{true} group, so a mis-assignment is penalised rather than hidden.

\textbf{Supervision required.} The attribute on the \emph{training} split only --- not on the
calibration split, never at test. The reduction is from ``attribute labels on calibration
\emph{and} test data'' to ``attribute labels on training data only'': a meaningful weakening of the
requirement, not its removal.

\textbf{Probe error.} It propagates asymmetrically: the minority stratum carries the \emph{largest}
threshold, so a minority point misplaced into a majority stratum is judged against a tighter
threshold and under-covered, while the reverse error only wastes efficiency. We checked whether
probe accuracy predicts the penalty and it does \textbf{not} --- the cell with the lowest AUROC
($0.945$) has no failures, while a cell at AUROC $0.999$ has one. What travels with the penalty is
the arm's own group-dependence: the two largest failures are ERM. We say so rather than offer AUROC
as a sufficient diagnostic.

\textbf{Privacy and ethics.} Written as a constraint on use rather than a disclaimer. Two points are
stated plainly: CelebA's \emph{Male} field is a third-party binary annotation and \textbf{not
self-identified gender}, so a study that treats it as gender has erred before any ethical question
arises; and inferring such an attribute at test time assigns a sensitive label to someone who has
not provided it, wrongly for $2$--$5\%$ of people even at the AUROC we measure. The recommendation
is scoped to settings where the attribute is already lawfully held, or where the inference stays
inside an audit an operator runs on their own system, and it is explicitly ruled out where the
inferred label would be stored, disclosed, or used to route individual outcomes. We note that
special-category data regimes may prohibit the inference outright regardless of accuracy.

\rpoint{R2.6}{moderate the claims; improve reproducibility}

\begin{reviewer}
``Claims such as `training buys efficiency' and `accuracy is a good proxy' should be presented as
empirical tendencies \ldots Please also provide the exact calibration/test construction, per-group
counts, training hyperparameters, the shift-robust procedure, random seeds, and a link to archival
code with reproducible commands.''
\end{reviewer}

\where{Sections~5.4, 5.6 and 5.7 (moderation); Appendix~D (reproducibility).}

\textbf{Moderation.} Four statements are now withdrawn rather than softened, because measurement
contradicts them.

\begin{enumerate}[leftmargin=1.4em,itemsep=0.35em]
\item The per-cell accuracy--set-size correlation, reported bare in the submitted version, is
\textbf{uninformative in five of eight cells} once a cluster-bootstrap CI is attached --- each rests
on only three to five methods. We now report a \emph{direction}, not an estimate.
\item ``The absolute efficiency range is narrow'' is \textbf{withdrawn}. The spread runs $0.036$ to
$0.364$, and a two-one-sided-test equivalence check rejects equivalence at a $0.10$ margin in six of
eight cells.
\item ``Training buys efficiency, not coverage'' is now stated as \textbf{conditional on Mondrian}:
under marginal calibration training still moves coverage, so it is not a general claim.
\item ``Mondrian is $2$--$4\times$ more split-stable than marginal calibration'' is
\textbf{withdrawn}. It held on the two backbones of the submitted version but not on four
(Section~5.7).
\end{enumerate}

\begin{revised}
An earlier version also claimed Mondrian was $2$--$4\times$ more split-stable than marginal
calibration. On the four-backbone data the two are comparable --- marginal's SD is
$0.009$--$0.042$ --- and in two of the eight cells Mondrian is the less stable of the two, so we
withdraw that comparison. Mondrian's advantage here is in the \emph{level} it attains, not in its
variance across splits.
\end{revised}

That last one was found only because the reviewers asked for more backbones (R1.3), which is a fair
illustration of why the request mattered. ``Accuracy is a good proxy for conformal burden'' is
likewise a tendency: point-estimate inversions occur in five of eight cells but clear their
cluster-bootstrap interval in only one, by $0.014$ of coverage.

\textbf{Reproducibility.} Appendix~D now gives the exact calibration/test construction --- disjoint
reservoirs from a seeded permutation, resampled within their own group indices, truncated to a
common size, realised $\rho$ logged per draw, separate uniform streams for the randomised scores ---
the per-group counts for both datasets, every head hyperparameter including the fine-tuning
optimisers, the shift-robust threshold in closed form, the seed sets, and the measured run times. We
also report that the worst group's Mondrian calibration count never falls below $50$ (minimum $66$
on Waterbirds, $373$ on CelebA), so the pooled-quantile fallback is never triggered by these runs.

\textbf{Archival code.} A DOI will be minted from a tagged release and inserted before publication;
the manuscript currently carries a visible placeholder rather than a provisional link.

\section*{Reviewer 3}

\rpoint{R3.1}{the calibration-versus-representation claim is too strong; provide CIs or equivalence
tests}

\begin{reviewer}
``The claim that coverage is a calibration property rather than a representation property seems too
strong \ldots Mondrian also does not always reach the target coverage. The authors should narrow the
claim and provide stronger statistical tests, such as confidence intervals or equivalence tests.''
\end{reviewer}

We agree on all three counts.

\where{Title, Section~5.1 heading and fifteen further places (narrowing); all main tables (CIs);
Section~5.6 (equivalence test); Section~5.1 and Appendix~B (sub-target level).}

\textbf{Narrowed.} The title is now comparative, and the same softening was applied throughout
rather than to the title alone. Section~5.1's heading, for instance, has changed from ``Coverage Is
a Calibration Property, Not a Representation One'' to ``Calibration Governs Worst-Group Coverage Far
More than Training Does''.

\textbf{Confidence intervals.} Every point estimate in the main tables now carries a two-stage
cluster bootstrap interval that resamples training seeds and then calibration splits within them.

\textbf{Equivalence test.} Added: a two-one-sided-test check of whether the across-method spread in
worst-group set size lies inside a stated margin. It \textbf{rejects} equivalence at $0.10$ in six
of eight cells, which is why we withdrew the description of the efficiency range as narrow.

\textbf{Mondrian not reaching target.} Addressed as under R1.1: the mean over groups is $0.9024$
against a nominal $0.900$, so the sub-target \emph{minimum} is a min-over-groups selection effect.
We continue to report the minimum, because it is the operationally relevant number, but now always
alongside the mean so the two are not confused.

\rpoint{R3.2}{``training buys efficiency'' is specific to Mondrian; the correlation is weak in one
setting}

\begin{reviewer}
``The claim that training buys efficiency rather than coverage also seems specific to Mondrian
calibration \ldots Moreover, the accuracy--set size correlation is weak in one of the four settings.
The authors should state this claim as conditional on Mondrian calibration and provide statistical
significance or uncertainty for the reported correlations.''
\end{reviewer}

Both implemented, and the reviewer's concern proved larger than stated.

\where{Abstract, introduction, Section~5.4 and the discussion (conditioning); Table~5
(intervals).}

The claim is now \textbf{explicitly conditional on Mondrian}. And with cluster-bootstrap intervals,
the correlation is uninformative not in one of four settings but in \textbf{five of eight} ---
including the $-0.97$ and $-1.00$ we previously highlighted, which at $n=4$ methods are not evidence
of a deterministic relationship. Table~5 now prints each interval, italicised where it spans
essentially $[-1, +1]$, and the text reports a direction rather than an estimate.

\section*{Changes we made without being asked}

Three, all disclosed so that no reader discovers a silent alteration.

\begin{enumerate}[leftmargin=1.4em,itemsep=0.35em]
\item \textbf{A definitional inconsistency in Section~3.} The coverage gap was defined in a way
that contradicted both the implementation and Appendix~B. It now reads
$(1-\alpha) - \min_g \mathrm{cov}_g$, matching both.
\item \textbf{A rendering defect in the predicted-group table.} It declared eight columns and
supplied seven, leaving a stray empty column. Fixed during the rewrite; all nine tables now pass a
column-count check.
\item \textbf{An empty-set disclosure.} On DINOv2/Waterbirds the mean worst-group set size is
\emph{below one} ($0.926$--$0.964$ across scores), meaning a small fraction of prediction sets are
empty. This is a legitimate outcome of standard split conformal with a highly accurate model, and
coverage there remains valid, but a reader expecting sets of size $\ge 1$ would otherwise be
confused. We report it rather than clipping sets to a minimum of one, which would hide it.
\end{enumerate}

\vspace{0.4em}
\noindent\textcolor{rulegrey}{\rule{\textwidth}{1pt}}

We are grateful for the care the reviewers took. The revision is materially more honest than the
submission: two analyses were corrected, four claims withdrawn, and one pre-specified criterion
reported as failing in seven of eight cells. We believe the paper's central finding is stronger for
it, since it now rests on eight backbone--dataset cells and on a study in which the representation
itself is retrained, rather than on the absolute dichotomy the reviewers rightly questioned.

\end{document}
